\section{Methodology}\label{si:methodology}
Continuous pointwise state constraints can produce measure-valued
multipliers \citep{casas1986control}. Regularization and discretization
provide approximations \citep{meyer2006optimal,ito2003semi,deckelnick2009convergence}.
The calculations below impose bounds on the discrete state and evaluate
the corresponding discrete Karush--Kuhn--Tucker (KKT) conditions.

\subsection{Models and discretization}
\input{sections/cht_mesh_methods}

\subsubsection{Cartesian trajectories and independent optimization}
\label{si:transient}
The main formulation orders unknowns by time level and then by the
spatial C-order index. The last spatial coordinate varies fastest.
Let $N_s=n_s^3$ and let $S$ be the lower shift matrix on the
$N_t$ time levels. With $D_t=\diag(\Delta t_n)$, the temporal
difference matrix is $L=D_t^{-1}(I-S)$.
The prescribed initial value contributes only to the first right-hand
side block, $f^1=Cy^0/\Delta t_1$. These definitions give the
assembled operator and its independent slabwise actions in main
Equations~\Mref{eq:cht_backward_euler}--\Mref{eq:cht_trajectory_operator}.
Forward substitution verifies the recovered control:
\[
 y^n=(A_s+C/\Delta t_n)^{-1}
       \left(u^n+Cy^{n-1}/\Delta t_n\right).
\]
The transpose action starts from the final adjoint block and uses
$-Cp^{n+1}/\Delta t_{n+1}$ in the preceding block.

For a complete trajectory, the objective's tracking and control
weights are $W=(D_t/\overline{\Delta t})\otimes I_{N_s}$.
The algebraic adjoint satisfies $p=\alpha WU$. With unequal steps,
this convention distinguishes the multiplier of the assembled
unweighted state equation from a time-weighted adjoint variable.
The physical objective restores the factor
$h^3\overline{\Delta t}$. The state constraints apply at the
$N_t$ positive time levels. The initial value is prescribed
independently.

The independent reference solves the bounded least-squares problem
\begin{equation}\label{eq:si_cht_bvls}
 \min_{Y\leq\psi}\frac12\left\|
 \begin{bmatrix}W^{1/2}\\\sqrt\alpha W^{1/2}\mathcal A\end{bmatrix}Y
 -
 \begin{bmatrix}W^{1/2}d\\\sqrt\alpha W^{1/2}f\end{bmatrix}
 \right\|_2^2.
\end{equation}
The SciPy bounded-variable least-squares method uses tolerance
$10^{-13}$ and an iteration cap of 1000. Its active-set indicator
defines the reference multiplier: $\lambda_i=(\ell-MY)_i$ at an
upper-bound index identified by the reference solver and zero elsewhere. The five KKT
components are evaluated independently for this reference solution.
Three iterative methods solve the complete PDAS problem for comparison:
Jacobi-preconditioned conjugate gradient (CG), denoted Jacobi-CG, RefDef and
recycling.

The verification study contains three steady $4^3$ problems,
three uniform-step $3^3$-by-4-slab problems, and three unequal-step
problems on the same trajectory grid. The queries are
$\theta=0,1/2,1$. The transient initial temperature is
$2\times10^{-4}$, the positive-level bound is $10^{-4}$, and the
unequal steps are $(0.01,0.02,0.03,0.04)$. Uniform steps have length
$0.025$. The reference ranks are 10 for steady problems and 12
for trajectories. These are nine distinct optimization
problems and 27 PDAS solver runs.


\subsubsection{Cartesian targets and fixed physical bounds}
\label{si:targets}
The initial Gaussian centers are
$(0.3,0.3,0.3)$, $(0.3,0.7,0.3)$, $(0.7,0.3,0.7)$ and
$(0.7,0.7,0.7)$. At query $\theta$, a rotation of $1.2\theta$
radians about $(x_1,x_2)=(1/2,1/2)$ changes their first two
coordinates, and $0.035(2\theta-1)$ shifts their third coordinates.
Their weights are $(1,0.8+0.15\theta,0.65,0.45)$ and their common
Gaussian width is $\sigma=0.12$:
\[
 d(x,\theta)=\sum_{j=1}^4 a_j(\theta)
 \exp\left(-\frac{\|x-z_j(\theta)\|_2^2}{2\sigma^2}\right).
\]
Unequal weights and the vertical shift distinguish the 16
query targets. The \href{https://github.com/teeratornk/deflation-example/blob/ba682dd05b168d2f43ff0c58050813deee785223/docs/supplementary-figures.md#cartesian-target-fields}
{repository figures} show desired temperatures and active sets for three
steady queries and three time levels of one trajectory.

For a transient query, let $\tau=t/T_f$.
An additional rotation of $0.8\tau$ changes the first two center
coordinates. A shift $0.05\sin(2\pi\tau)$ changes the third
coordinate, and $0.15+0.85\sin^2(\pi\tau)$ multiplies every
Gaussian weight. One query specifies this complete moving and
pulsed trajectory. The comparison holds $T_f=0.1$ fixed.

A direct-PDAS bisection with eight steps calibrates a fixed bound
near active fraction 0.2 at $\theta=1/2$. Calibration uses $24^3$
nodes for steady problems and $12^3$ nodes with eight slabs for
trajectories. Each refinement uses the corresponding bound unchanged.
Calibration precedes the solver comparisons; its measured cost enters
the preparation-inclusive times once for each method.

\subsection{Reference construction and analysis}
\input{sections/cht_mesh_reference_methods}

\subsubsection{Reference eigenpairs and compact storage}
\label{si:construction}
Let $L_s\phi_j=\mu_j\phi_j$ with orthonormal spatial sine modes.
For any temporal vector $q$,
\[
 \mathcal A_0(q\otimes\phi_j)
 =((\mu_jI+c_0L)q)\otimes\phi_j.
\]
Applying the weighted normal operator gives
$M_0(q\otimes\phi_j)=(H_jq)\otimes\phi_j$, with $H_j$ defined
in main Equation~\Mref{eq:cht_temporal_eigenproblem}.
Thus the mode-dependent construction diagonalizes the specified
homogeneous space--time reference exactly, up to floating-point
eigensolver error.

The candidate cutoff follows from monotonicity. Since
$W_t=D_t/\overline{\Delta t}$,
\[
 W_tL+L^\top W_t
 =\frac{2I-S-S^\top}{\overline{\Delta t}}
\]
is positive definite. For $\mu\geq0$,
\[
 \frac{\partial}{\partial\mu}
 \left[W_t+\alpha(\mu I+c_0L)^\top W_t(\mu I+c_0L)\right]
 =\alpha\left(2\mu W_t+c_0(W_tL+L^\top W_t)\right)
\]
is positive definite. The minimum eigenvalue therefore increases
with the spatial Laplacian eigenvalue. A spatial mode beyond the
first $r$ has at least $r$ preceding temporal candidates with
no larger reference eigenvalue. The first $\min(r,N_s)$ spatial
modes consequently suffice to identify a lowest rank-$r$ reference
space. A fixed lexicographic order resolves repeated spatial
eigenvalues at the cutoff.

The compact reference stores each selected temporal vector $q_{j,k}$
with its spatial factor $\phi_j$. At spatial node $a$ and time level $n$,
restriction evaluates $\phi_j(a)q_{j,k}(n)$ directly. Only the current
inactive candidate matrix is assembled. Memory measurements include
these factors, the restricted candidate, the orthonormal basis and the
cached operator product.

\subsubsection{Complete and truncated temporal coverage}
\label{si:temporal_coverage}
Let $\Phi_S=[\phi_1,\ldots,\phi_S]$ collect the retained spatial modes and let
$Q_j\in\mathbb R^{N_t\times N_t}$ hold all $N_t$ temporal factors of mode $j$ for a
trajectory with $N_t$ time levels. The factors are orthonormal in the temporal
metric, so every $Q_j$ is invertible. The complete reference therefore spans
\[
 \operatorname{span}\{q_{j,k}\otimes\phi_j\}
 =\operatorname{span}\{e_n\otimes\phi_j:\ n=1,\ldots,N_t,\ j=1,\ldots,S\},
\]
the space that pairs every retained spatial mode with every time level. A
restriction $R$ selects rows, and
$R\,[\,Q_1\otimes\phi_1,\ldots,Q_S\otimes\phi_S\,]
=R\,[\,I\otimes\phi_1,\ldots,I\otimes\phi_S\,]\operatorname{blkdiag}(Q_j\otimes 1)$
with an invertible right factor. The restricted bases span the same space, so every
coarse correction built from them, $Z(Z^\top BZ)^{-1}Z^\top$ and the projector $P$
of main Equation~\Mref{eq:revision_projectors}, is identical in exact arithmetic.
The complete construction thus depends only on the retained spatial modes; with
complete temporal coverage, the temporal eigenvectors provide an alternative basis of the
same correction space.

A truncated construction keeps fewer than $N_t$ factors per mode, or selects the
lowest spatial--temporal pairs up to a total rank. The temporal eigenvalue ordering
then decides which time dependence of each mode survives, and the span no longer
contains every $e_n\otimes\phi_j$. The companion tests confirm both statements:
on random inactive sets of a transformer trajectory, the restricted complete and
tensor bases give coarse corrections that agree to rounding, and a truncated
selection misses time-level indicators of a retained mode. The Cartesian
trajectories with $S$ spatial modes and the Bore~3 trajectory use complete
coverage; the transformer, Bore~1 and Bore~2 trajectories select the 100 lowest
pairs.

\subsubsection{Restriction analysis and conditional convergence}
\label{si:theory}
Let $M_0,M\in\mathbb R^{N\times N}$ be symmetric positive definite (SPD)
and set $\Delta=M-M_0$.
Let $\Phi^\top\Phi=I_r$ and
$G_0=M_0\Phi-\Phi\Lambda$ for symmetric $\Lambda$.
For the inactive and active row partitions, main
Equation~\Mref{eq:restriction_identity} follows directly from
block multiplication. If $\Phi_I=ZR$ has full column rank, the
subspace residual satisfies
\begin{equation}\label{eq:si_F_bound}
 \|BZ-Z(Z^\top BZ)\|_F
 \leq\frac{
 \|(M_0)_{IJ}\Phi_J\|_F+
 \|\Delta_{II}\Phi_I\|_F+
 \|(G_0)_I\|_F}{\sigma_{\min}(\Phi_I)}.
\end{equation}
Projection onto the orthogonal complement of $Z$ eliminates
$\Phi_I\Lambda$. The factor $R^{-1}$ has norm
$1/\sigma_{\min}(\Phi_I)$, which gives the displayed estimate.
Partial cancellation between its three terms can make the
actual residual smaller than this sufficient bound.

Orthonormality of the full reference also gives
\begin{equation}\label{eq:si_massloss}
 \Phi_I^\top\Phi_I=I_r-\Phi_J^\top\Phi_J.
\end{equation}
For $\beta_J=\|\Phi_J\|_2<1$, the restriction has full rank and
$\sigma_{\min}(\Phi_I)=\sqrt{1-\beta_J^2}$. This condition describes
the least retained norm of any unit combination of modes. Individual
column norms alone do not determine the rank of their combined
restriction.

Let $U_r$ span the lowest $r$ eigendirections of $B$, with
$1\leq r<|I|$, and let $U_\perp$ span its orthogonal complement.
Write $E=Z^\top BZ$, $F=BZ-ZE$ and
$BU_\perp=U_\perp\Lambda_\perp$. Premultiplication yields
\begin{equation}\label{eq:si_sylvester}
 \Lambda_\perp X-XE=U_\perp^\top F,\qquad X=U_\perp^\top Z.
\end{equation}
If $\eta=\lambda_{r+1}(B)-\lambda_{\max}(E)>0$, diagonalizing $E$
divides each transformed entry on the right by an eigenvalue
difference at least $\eta$. The Frobenius norm then gives
\[
 \|\sin\Theta(U_r,Z)\|_F=\|X\|_F\leq\|F\|_F/\eta.
\]
This is a sufficient estimate for the specified low invariant
subspace. Rank loss requires a lower-rank comparison. A repeated
eigenvalue at the cutoff requires a choice of target subspace or
a complete-cluster comparison.

Combining these results bounds the largest angle $\vartheta$ by
\begin{equation}\label{eq:si_combined_angle}
 \sin\vartheta\leq\varepsilon_Z,\qquad
 \varepsilon_Z=\min\left\{1,
 \frac{\|(M_0)_{IJ}\Phi_J\|_F+
       \|\Delta_{II}\Phi_I\|_F+\|(G_0)_I\|_F}
      {\sigma_{\min}(\Phi_I)\eta}\right\},
 \quad \eta>0.
\end{equation}
A bound equal to one is trivial. A nonpositive separation leaves
this particular estimate unavailable.
For deflated CG, the perturbation estimate of
\citet{kahl2017deflated} gives
\begin{equation}\label{eq:si_effective_kappa}
 \kappa_{\rm eff}\leq
 \left(\sqrt{\frac{\lambda_{\max}(B)}{\lambda_{r+1}(B)}}+
       \sqrt{\kappa_2(B)}\,\varepsilon_Z\right)^2.
\end{equation}
Here $\kappa_{\rm eff}$ is the ratio of the largest to smallest
nonzero deflated eigenvalues. Following the coarse correction,
the usual exact-arithmetic CG estimate is
\[
 \|e_k\|_B\leq2
 \left(\frac{\sqrt{\kappa_{\rm eff}}-1}
            {\sqrt{\kappa_{\rm eff}}+1}\right)^k\|e_0\|_B.
\]
A conditioning-improvement certificate requires the upper bound
in \eqref{eq:si_effective_kappa} to be smaller than the original
condition number. A nontrivial angle bound alone does not imply
that comparison.

For Jacobi-preconditioned CG, every matrix, subspace residual,
angle and separation in this test uses
$\widehat B=D^{-1/2}BD^{-1/2}$ and an orthonormal basis of
$D^{1/2}Z$. The rank-one controlled CHT example starts with an
exact lowest eigenvector of the full Jacobi-scaled CHT matrix.
It removes nodes in increasing order of the absolute reference
amplitude and applies \eqref{eq:si_combined_angle}, including
the numerical reference residual. Its physical coarse vector is
obtained by the inverse Jacobi square-root scaling. The analytical
timing references have their separately stated construction.

\subsubsection{Projected recurrence and numerical safeguards}
\label{si:implementation}
For $E=Z^\top BZ$, $Q=ZE^{-1}Z^\top$ and $P=I-QB$,
$QBQ=Q$, $P^2=P$ and $BP=P^\top B$.
The balanced inverse $P\mathcal JP^\top+Q$ is SPD for SPD $\mathcal J$:
both terms have nonnegative quadratic forms, and their nullspaces
intersect only at zero. If $Z^\top r=0$, then $P^\top r=r$ and
$Qr=0$, giving the projected action $P\mathcal Jr$ used in the implementation
\citep{tang2009comparison,gaul2014framework}.
The corrected initial guess and coarse-corrected restarts supply
this residual condition.

A thin QR factorization into orthonormal and upper-triangular factors,
followed by a small-factor singular value decomposition (SVD),
determines numerical rank on the GPU.
RefDef and recycling apply the same relative threshold
$10^{-12}$ and coarse condition limit $10^{10}$.
Cholesky solves evaluate the coarse correction. Retaining $BZ$
changes the implementation of a matrix product in the projector;
its setup and storage are charged for both methods. The efficient
rank-zero path allocates no coarse basis or cached coarse product.
Original-residual checks apply to every returned inactive solution.
PDAS then verifies the five KKT components in main
Equation~\Mref{eq:cht_kkt_acceptance}.

The recycling candidate pool retains the deployed coarse vectors
and the last $r$ newly generated search directions. Selection uses
the current Jacobi-scaled operator. The selected physical vectors
are transferred by zero extension between consecutive PDAS inner
solves. The implementation keeps the coarse space fixed during
an inner solve and its restarts. A failed target clears the
solution and recycling history before the next query.

\subsection{Implementation and measurement}
\subsubsection{Overview of the comparisons}
Table~\ref{si:tab_experiment_map} lists the problems, references and measurements of
the principal comparisons; the Results section describes the additional studies.

\begin{table}[pos=htbp]
\centering\small
\caption{Comparisons, problems and measurements. Every timed comparison
uses 16 targets and five complete sequence repetitions per method unless
stated otherwise. Ranks denote total spatial or space--time rank, and $S$
the number of spatial modes, each paired with a temporal factor for every slab.}
\label{si:tab_experiment_map}
\begin{tabularx}{\textwidth}{@{}>{\raggedright\arraybackslash}p{0.19\textwidth}>{\raggedright\arraybackslash}p{0.25\textwidth}>{\raggedright\arraybackslash}X@{}}
\toprule
Purpose & Problems and reference & Measurements \\
\midrule
Optimization verification &
12 body-fitted and nine Cartesian steady or transient verification problems &
Independent bounded least squares, original inner residuals,
KKT components, state and objective agreement, trajectory recovery \\
Body-fitted optimization &
Axisymmetric transformer and three-dimensional bore-in-block meshes; coarse P1 references &
Complete steady and transient costs, KKT and original-residual checks,
constraint changes, rank effects, and process memory \\
Steady scale &
Cartesian $32^3$--$128^3$; transformer; Bore~1--4 up to
2.5 million DOFs; ranks 100 and 200 &
Complete time and inner iterations under main Algorithm~\Mref{alg:final_procedure},
original residuals and KKT checks \\
Trajectories &
Cartesian $12^3$--$64^3$ with 8--32 slabs, $S=12$ or 6;
transformer and Bore~1--3 with 4 or 8 slabs &
Complete trajectory optimization, original residuals, KKT and recovery checks \\
Operating conditions &
Transformer day, drive cycle, off-design flow, overloads and query order;
one nominal reference &
Complete time over chained windows or reassembled operators;
original residuals and KKT checks \\
Retention ablation &
Six primary body-fitted sequences; the same reference by direct restriction or
zero-extension transfer, with Jacobi-CG in the same job &
Complete time and inner iterations \\
Transfer mechanism &
Steady and transient PDAS traces; identical replay systems &
Activation and release, retained rank, coarse energy removal,
deflated spectrum on small systems, original-residual acceptance \\
Conditional analysis &
Controlled $6^3$ CHT restrictions; exact rank-one scaled reference &
Separation, actual angle, sufficient bound and effective conditioning \\
\bottomrule
\end{tabularx}
\end{table}

\subsubsection{Numerical environment and AmgX configuration}
\label{si:cht_environment}
The complete performance comparisons use one NVIDIA H200 GPU and
four threads on an Intel Xeon Platinum 8468 CPU. The software
environment uses Python 3.12.3, NumPy 2.5.3, SciPy 1.18.1,
PyTorch 2.10.0 with CUDA 12.8 and AmgX 2.5.0. NumPy and SciPy use
OpenBLAS libraries reporting versions 0.3.34.106.0 and 0.3.31.dev.
The body-fitted comparisons use the same hardware.
The \href{https://github.com/teeratornk/deflation-example/blob/ba682dd05b168d2f43ff0c58050813deee785223/docs/mesh-showcases.md#reproduce-the-current-comparisons}
{body-fitted reproduction examples} give the commands of the timed comparisons; the
\href{https://github.com/teeratornk/deflation-example/blob/ba682dd05b168d2f43ff0c58050813deee785223/docs/mesh-showcases.md#reproduce-the-measured-implementations}
{measured implementations} of the earlier campaign and of the finer-transformer
residual-correction study are listed separately. A separate
\href{https://github.com/teeratornk/deflation-example/blob/ba682dd05b168d2f43ff0c58050813deee785223/docs/mesh-showcases.md#use-the-recommended-guarded-solver}
{guarded-solver example} checks the initial guess before solving.
The \href{https://github.com/teeratornk/deflation-example/blob/ba682dd05b168d2f43ff0c58050813deee785223/docs/reference-policy-study.md#frozen-complete-optimization-campaign}
{Cartesian examples} reproduce the spatial, temporal and rank comparisons.
These examples provide the corresponding configurations, data and numerical procedures.

Table~\ref{si:tab_amgx_configuration} specifies the native solver.
The session retains its base configuration and resource handles.
Each inner solve creates a configuration with its right-hand-side
dependent absolute target, uploads the current matrix and initial
guess, constructs a hierarchy and solves. Matrix, vector, solver
and per-solve configuration handles are destroyed after verification.
The session closes after the last target. Every method process
initializes and warms both GPU libraries before its complete
sequence interval.

\begin{table}[pos=htbp]
\centering\small
\caption{The complete steady and transient comparisons use the
same AmgX hierarchy settings. Original inactive systems use the absolute
target $10^{-11}\|b\|_2$, with $10^{-11}$ used when $b=0$.
The separate correction protocol in Section~\ref{si:mesh_refinement_diagnosis}
specifies error-equation targets and remaining iteration budgets.}
\label{si:tab_amgx_configuration}
\begin{tabularx}{\textwidth}{@{}lX@{}}
\toprule
Setting & Value \\
\midrule
Deterministic mode & Enabled \\
Outer solver; arithmetic & Preconditioned conjugate gradient (PCG); double precision \\
Convergence; residual norm; iteration cap &
Absolute; Euclidean; 10000 for Cartesian, 50000 for body-fitted \\
Preconditioner; construction & Algebraic multigrid (AMG); classical \\
Coarsening; interpolation & Parallel maximal independent set (PMIS); distance-two (D2) \\
Smoother; relaxation & Symmetric multicolor Gauss--Seidel; 1 \\
Pre-sweeps; post-sweeps & 1; 1 \\
Cycle; cycles per application; maximum levels & V; 1; 25 \\
Coarse solver & Dense lower--upper (LU) factorization \\
\bottomrule
\end{tabularx}
\end{table}

\Needspace{9\baselineskip}
\subsubsection{Disjoint timing intervals and memory measurements}
\label{si:cht_timing}
Table~\ref{si:tab_timing_definitions} partitions complete-sequence
time.
The outer timers exclude the enclosed linear-solver call; the
kernel timers partition that call. A synchronization barrier
immediately precedes each GPU phase. Explicit waiting belongs to
the synchronization interval. Waiting inside a numerical-library
call belongs to that call's interval.

\begin{table}[pos=H]
\centering\small
\caption{The timing intervals partition the measured complete
optimization sequence. Their sum equals the complete time.
Preparation costs use separate enclosing process measurements.}
\label{si:tab_timing_definitions}
\begin{tabularx}{\textwidth}{@{}p{0.28\textwidth}X@{}}
\toprule
Interval & Included work \\
\midrule
Assembly and target construction &
Spatial or space--time operators, reduced matrix, desired states
and right-hand sides \\
Reference construction &
Full-domain spatial and temporal factors \\
Inactive-system and coarse-space setup &
Inactive matrices and right-hand sides; reference restriction or
recycling transfer and selection; rank-revealing orthogonalization,
coarse products and factorization, or AMG hierarchy construction \\
Initialization and iteration &
Initial correction and residual; iterative solver recurrence \\
Conversion and transfers &
Matrix and vector conversion; uploads and solution downloads \\
Verification and active-set updates &
Independently recomputed residuals, state and multiplier recovery,
KKT components and constraint updates \\
Remaining work &
Solver resource creation, synchronization, cleanup and other host work \\
\bottomrule
\end{tabularx}
\end{table}

Component figures use the actual repetition with median complete
time. Thus their stacked intervals add to the displayed total.
Separate component medians summarize their own distributions.
The complete-time ranges include all five repetitions.

The speedup of a comparison is $\operatorname{med}T_{\rm alt}/\operatorname{med}T_{\rm ref}$,
the ratio of the median complete times of the fastest alternative and RefDef over their converged repetitions. The bracketed interval is the envelope
\[
 \left[\frac{\min T_{\rm alt}}{\max T_{\rm ref}},\;
        \frac{\max T_{\rm alt}}{\min T_{\rm ref}}\right]
\]
over the same repetitions. It combines extreme times of separate repetitions; it is
an envelope of observed ratios, without pairing or a confidence level. The timing
ranges overlap when the envelope contains one. Comparisons with a single repetition per method
report the ratio alone.

The enclosing process interval runs from launch to termination.
The difference between this interval and its explicitly timed
initialization, warmup, complete sequence and finalization includes
imports, process creation and monitoring setup. Adding the common
calibration once gives the preparation-inclusive total,
$T_{\rm prep}=T_{\rm launch\text{-}to\text{-}exit}+T_{\rm cal}$. The common
calibration belongs to the family and is charged equally to every
method in a single-deployment comparison.

A separate monitor samples host resident memory and total GPU process
allocation at a nominal period of 10~ms. The largest sampling gap
quantifies the temporal resolution of each measurement. GPU allocation includes allocator
caches and all numerical backends in the process. Host and GPU
peaks are sampled lower bounds. The same measurement and accounting
boundary applies to all four methods.
